\documentclass[11pt]{amsart}
\usepackage{amssymb,amsmath,amsthm}
\usepackage{mathtools}
\usepackage{booktabs}
\usepackage[margin=1.15in]{geometry}
\usepackage[colorlinks=true,linkcolor=blue,citecolor=blue,urlcolor=blue]{hyperref}

\newtheorem{theorem}{Theorem}[section]
\newtheorem{lemma}[theorem]{Lemma}
\newtheorem{proposition}[theorem]{Proposition}
\newtheorem{corollary}[theorem]{Corollary}
\theoremstyle{definition}
\newtheorem{definition}[theorem]{Definition}
\theoremstyle{remark}
\newtheorem{remark}[theorem]{Remark}
\numberwithin{equation}{section}
\raggedbottom

\newcommand{\N}{\mathbb{N}}
\newcommand{\R}{\mathbb{R}}
\newcommand{\cA}{\mathcal{A}}
\newcommand{\cJ}{\mathcal{J}}
\newcommand{\cL}{\mathcal{L}}
\newcommand{\cB}{\mathcal{B}}
\newcommand{\dd}{\,\mathrm{d}}
\newcommand{\Leb}{\mathrm{Leb}}
\DeclareMathOperator{\supp}{supp}

\title[Hankel determinants with power-law atoms]{Hankel determinants of discrete measures with power-law atoms, and of zeta values}
\author{Gordon [surname] (draft)}
\date{Draft of 14 October 2026}



\begin{document}
\begin{abstract}
Let $\sigma=\sum_{k\ge1}w_k\delta_{x_k}$ be a discrete measure whose atoms are regularly spaced with
$x_k\sim Ck^{-p}$ for some $p>0$, and whose weights satisfy $\sum_k w_k<\infty$ and $w_k\ge c_1k^{-s}$. We prove that
the Hankel determinants $D_n=\det(\int x^{i+j}\dd\sigma)_{i,j=0}^{n-1}$ satisfy
\[
  \log D_n=-pn^2\log n+\Bigl(\tfrac{3p}{2}+\log\bigl(c(p)C\bigr)\Bigr)n^2+o(n^2),\qquad
  c(p)=\tfrac14\Bigl(\tfrac{\sqrt\pi\,\Gamma(1/p+1)}{\Gamma(1/p+1/2)}\Bigr)^{p}.
\]
The constant comes from an explicitly solvable constrained equilibrium problem: the minimiser of the logarithmic
energy $\iint\log\frac{1}{|u^{-p}-v^{-p}|}\dd\mu(u)\dd\mu(v)$ among probability measures $\mu\le\dd u$ on $(0,\infty)$ is
unique, it is an average of arcsine laws, and its density is a regularised incomplete beta function. As a
corollary, for every real $a>0$ and real $b$ with $2a+b>1$,
$\log\det_{1\le i,j\le n}\zeta(a(i+j)+b)=-an^2\log n+(\tfrac{3a}{2}+\log c(a))n^2+o(n^2)$. For positive integers $a,b$
this gives the leading-order behaviour $-an^2\log n$ expected by Haynes and Zudilin, together with the $n^2$ term.
\end{abstract}
\maketitle

\section{Introduction}\label{sec:intro}

\subsection{Setting and main result}
For a positive Borel measure $\sigma$ on $\R$ with finite moments $m_r=\int x^r\dd\sigma(x)$, write
\[
  D_n(\sigma):=\det\bigl(m_{i+j}\bigr)_{i,j=0}^{n-1}\qquad(n\ge1)
\]
for its Hankel determinants. We study discrete measures
\begin{equation}\label{eq:sigma}
  \sigma=\sum_{k\ge1}w_k\,\delta_{x_k},\qquad x_k=f(k),
\end{equation}
whose atoms accumulate at a single point with polynomial spacing, the model case being $x_k=Ck^{-p}$. Throughout,
$p>0$ and
\begin{equation}\label{eq:constants}
  I(p):=\frac{\sqrt\pi}{2}\,\frac{\Gamma(\frac1p+\frac12)}{\Gamma(\frac1p+1)},\qquad
  c(p):=\frac14\Bigl(\frac{\pi}{2I(p)}\Bigr)^{p},\qquad
  \kappa_p:=\frac{3p}{2}+\log c(p).
\end{equation}
For instance $I(1)=\pi/4$, $c(1)=\tfrac12$, and $I(2)=1$, $c(2)=\pi^2/16$.

\begin{theorem}\label{thm:main}
Let $p>0$ and $C>0$. Let $f\in C^1[1,\infty)$ satisfy
\begin{equation}\label{eq:Ha}
  \lim_{\xi\to\infty}\xi^{p+1}f'(\xi)=-pC,\qquad
  0<\inf_{\xi\ge1}\xi^{p+1}|f'(\xi)|\le\sup_{\xi\ge1}\xi^{p+1}|f'(\xi)|<\infty,
\end{equation}
and let the weights $w_k\ge0$ satisfy
\begin{equation}\label{eq:Hw}
  W:=\sum_{k\ge1}w_k<\infty\qquad\text{and}\qquad w_k\ge c_1k^{-s}\quad(k\ge1)
\end{equation}
for some constants $c_1>0$ and $s\in\R$. Then the measure $\sigma$ of \eqref{eq:sigma} satisfies
\begin{equation}\label{eq:main}
  \log D_n(\sigma)=-pn^2\log n+\bigl(\kappa_p+\log C\bigr)n^2+o(n^2)\qquad(n\to\infty).
\end{equation}
\end{theorem}

The upper bound in the second condition of \eqref{eq:Ha} already follows from the limit and the
continuity of $f'$; we state it for symmetry. Under \eqref{eq:Ha}, $f$ is strictly decreasing with a finite limit $L$, and $f(\xi)-L\sim C\xi^{-p}$; Hankel
determinants are invariant under translation of $\sigma$ (Lemma~\ref{lem:elem}), so $L$ plays no role. The upper bound
in \eqref{eq:main} uses only $\sum_kw_k<\infty$; apart from the finiteness of the moments, the lower bound uses only
$w_k\ge c_1k^{-s}$.

Since $\sigma$ has infinitely many atoms, it has monic orthogonal polynomials of all degrees and a Jacobi matrix with
off-diagonal entries $b_j>0$ (Lemma~\ref{lem:jacobi}). Theorem~\ref{thm:main} is equivalent to a doubly averaged
statement about these coefficients.

\begin{corollary}\label{cor:jacobi}
Under the hypotheses of Theorem~\ref{thm:main},
\[
  \frac{2}{n^2}\sum_{j=1}^{n-1}(n-j)\log\bigl(j^pb_j\bigr)\longrightarrow\log\bigl(c(p)C\bigr).
\]
\end{corollary}

The constant $\kappa_p$ is the value of a variational problem, which we solve explicitly. Let
\[
  \cA:=\bigl\{\mu\ \text{positive Borel measure on }(0,\infty):\ \mu\bigl((0,\infty)\bigr)=1,\ \mu\le\Leb\bigr\},
\]
where $\mu\le\Leb$ means $\mu(B)\le|B|$ for every Borel set $B$, and for $\mu\in\cA$ put
\begin{equation}\label{eq:J}
  \cJ_p(\mu):=\iint\log\bigl|u^{-p}-v^{-p}\bigr|\dd\mu(u)\dd\mu(v)\in[-\infty,\infty).
\end{equation}
For $\beta>0$ let $\omega_\beta$ be the arcsine probability measure on $[0,\beta]$, with density
$1/(\pi\sqrt{y(\beta-y)})$. For $V\in[0,1]$ and $\alpha,\gamma>0$ let
$\cB_V(\alpha,\gamma):=B(\alpha,\gamma)^{-1}\int_0^Vs^{\alpha-1}(1-s)^{\gamma-1}\dd s$ be the regularised incomplete beta
function.

\begin{theorem}\label{thm:equilibrium}
Let $p>0$ and $c=c(p)$. Let $\mu^*$ be the probability measure on $(0,\infty)$
\[
  \mu^*:=\int_0^1\omega_{4ct^{-p}}\dd t,\qquad\text{i.e.}\quad\mu^*(B)=\int_0^1\omega_{4ct^{-p}}(B)\dd t,
\]
and let $\mu^*_u$ be its image under $x\mapsto x^{-1/p}$. Then:
\begin{enumerate}
\item $\mu^*_u\in\cA$, with density
  $\rho^*(u)=\cB_{\min(1,\,u^{-p}/4c)}\bigl(\tfrac1p+\tfrac12,\tfrac12\bigr)$. In particular $\rho^*=1$ on
  $(0,u_s]$, where $u_s:=(4c)^{-1/p}=2I(p)/\pi$, and $0<\rho^*<1$ on $(u_s,\infty)$.
\item The function $U^*(x):=\int\log|x-y|\dd\mu^*(y)$ (the negative of the logarithmic potential of $\mu^*$) satisfies
  $U^*(x)=\log c+p$ for $0<x\le4c$, and
  $\ge\log c+p$ for all $x>0$.
\item $\cJ_p(\mu^*_u)=\kappa_p$, and $\cJ_p(\mu)<\kappa_p$ for every $\mu\in\cA$ with $\mu\neq\mu^*_u$. Thus
  $\sup_{\cA}\cJ_p=\kappa_p$, attained only at $\mu^*_u$.
\end{enumerate}
\end{theorem}

Applied to the measure $\sum_kk^{-(2a+b)}\delta_{k^{-a}}$, Theorem~\ref{thm:main} gives the asymptotics of Hankel
determinants of zeta values along arithmetic progressions.

\begin{corollary}\label{cor:zeta}
Let $a>0$ and $b$ be real numbers with $2a+b>1$. Then
\[
  \log\det_{1\le i,j\le n}\zeta\bigl(a(i+j)+b\bigr)=-an^2\log n+\Bigl(\frac{3a}{2}+\log c(a)\Bigr)n^2+o(n^2).
\]
In particular, $\log\det_{1\le i,j\le n}\zeta(i+j+b)=-n^2\log n+(\frac32-\log2)n^2+o(n^2)$ for every $b>-1$, and
$\log\det_{1\le i,j\le n}\zeta(2(i+j)+b)=-2n^2\log n+(3+\log\frac{\pi^2}{16})n^2+o(n^2)$ for every $b>-3$.
\end{corollary}

\subsection{Context}\label{sec:context}
\emph{Hankel determinants of zeta values.} Monien~\cite{Monien} studied the Hankel determinants
$H_n^{(r)}[F]=\det(F(i+j+r))$ of Dirichlet series $F$, gave a multiple-sum formula for them \cite[Thm.~2.1]{Monien},
and wrote for $F=\zeta$: ``We do not know how to prove the full asymptotic expansions but will describe a method
that lets one at least understand the weaker asymptotic form $\log H_n^{(0)}[\zeta]\sim\log H_n^{(1)}[\zeta]\sim
-n^2(\log(2n)-3/2)$.'' He also reports numerical experiments of D.~Zagier consistent with this form. Haynes and
Zudilin~\cite{HaynesZudilin} wrote: ``for positive integers $a$ and $b$, we expect that $\log\det_{1\le i,j\le n}
\zeta(a(i+j)+b)=-an^2\log n+O(n^2)$ as $n\to\infty$'', and proved the weaker bound ``$\log\det_{1\le i,j\le n}
\zeta(a(i+j)+b)\le-(a/2)n^2\log n+O(n^2)$''. They also write
that ``In [5] Monien developed these ideas further and rigorously justified the above asymptotics in the case
$r=0$'', their reference [5] being a personal communication \cite[\S1]{HaynesZudilin}. Coons and Sutherland
\cite[Thm.~1.3]{CoonsSutherland} proved, for every nondegenerate Dirichlet series $F(s)=\sum f(n)n^{-s}$ with
$f\ge0$ convergent for $\Re s\ge s_0$, that $H_n^{(r)}[F]>0$ for $n\ge1$, $r\ge s_0-2$, and that
$\log H_n^{(r)}[F]<-cn^2$ for some $c>0$ and all sufficiently large $n$ and $r$.

Corollary~\ref{cor:zeta} therefore proves the asymptotic formula expected in~\cite{HaynesZudilin}, for all positive
integers $a,b$ and more generally for real $a>0$ and $2a+b>1$, and it identifies the $n^2$ term. For $a=1$ the
constant $\frac32-\log2$ agrees with Monien's form $-n^2(\log 2n-\frac32)=-n^2\log n+(\frac32-\log2)n^2$, provided
the latter is read with an $o(n^2)$ error, as the expansion from Zagier's numerics quoted in~\cite{Monien}
indicates. (Read literally, Monien's ``$\sim$'' carries only the leading term.) For $a=2$ and $b\in\{-2,0,2\}$ the determinants are known in closed form, by results of
Baricz and \v{S}tampach~\cite{BariczStampach} on Hankel determinants of Rayleigh functions (the closed forms we use
were confirmed independently against direct determinants); their asymptotics are
consistent with Corollary~\ref{cor:zeta} (Remark~\ref{rem:bs}). Apart from these closed-form cases, we have not found in
the literature a published proof of the leading term $-an^2\log n$ (for $a=1$, $r=0$ see the remark
of~\cite{HaynesZudilin} quoted above), nor of the $n^2$ term $\frac{3a}{2}+\log c(a)$.

Hankel determinants of zeta-type moments also occur in recent irrationality work. Lai~\cite{Lai} writes that
``Fauzan's proof uses Hankel determinants, which were previously studied by Haynes--Zudilin [23] and Zudilin [55]''.
Fauzan's announcement~\cite{Fauzan} is an unrefereed Zenodo record; the full manuscript was not available to us. Its
record page describes ``normalized Hankel determinants'' together with ``a positive moment representation and a
logarithmic-energy estimate''. The determinants there appear to be built from a different moment sequence. There is
nonetheless a methodological overlap with the present paper (Hankel determinants of zeta-type moments controlled by
a logarithmic-energy argument), and we cannot exclude that the announced energy estimate is related to
Theorem~\ref{thm:equilibrium}. The present paper is not concerned with arithmetic applications, and none of its
proofs uses~\cite{Fauzan} or~\cite{Lai}.

\emph{The equilibrium problem.} Since $-\cJ_p(\mu)$ is the logarithmic energy of the image measure $g_*\mu$, the
maximisation of $\cJ_p$ over $\cA$ is a logarithmic energy minimisation with an upper constraint on the measure.
Constrained energy problems of this kind are the subject of a substantial literature. We have only consulted the
abstracts of the following works, and the summary below is restricted to what those abstracts say. According to the
abstract of~\cite{CDLO}, the constrained equilibrium problem was introduced by Rakhmanov~\cite{Rakhmanov} to describe
the asymptotic zero distribution of discrete orthogonal polynomials in a compact interval, and was extended by
Dragnev and Saff~\cite{DragnevSaff} and by Kuijlaars and Van Assche~\cite{KVA-PLMS}. The abstract of~\cite{DragnevSaff}
announces existence and uniqueness of the constrained minimiser, a characterisation theorem and a saturation result,
with explicit results for Krawtchouk polynomials. The abstract of~\cite{KVA-PLMS} states that for extremal polynomials
on infinite discrete sets such as those of the Meixner and Charlier polynomials, the asymptotics are governed by a
constrained extremal energy problem that can be solved explicitly. The abstract of Van Assche and Van
Baelen~\cite{VAVB} concerns a $q$-lattice $\{q^k\}$, which accumulates at $0$. We use no result from this theory:
existence, uniqueness and the value of the extremum are proved directly in Section~\ref{sec:equilibrium}. We have not
found the problem $\sup_\cA\cJ_p$, whose constraint has infinite total mass near the accumulation point, or its
solution, in the literature. Because the full texts of \cite{Rakhmanov,DragnevSaff,KVA-PLMS,KVA-JAT} have not been
read, this statement is provisional.

For $p=1$, Monien~\cite[\S3]{Monien} set up a constrained maximisation of this type heuristically: ``We seek a
continuous test function $\rho(x)$ which maximizes the functional Eq.~3.5 subject to the constraint that $0<\rho(x)\le
1$ for all $x$ and $\rho(x)=0$ for $x<0$ and Eq.~3.4. We assume that such a continuous functions exists.'' His
Theorem~3.1 gives the solution of the stationarity condition,
\[
  \rho(x)=1\ \ (0\le x\le\tfrac12),\qquad
  \rho(x)=\tfrac{2}{\pi}\Bigl(\arctan\tfrac{1}{\sqrt{2x-1}}-\tfrac{\sqrt{2x-1}}{2x}\Bigr)\ \ (x>\tfrac12),
\]
which coincides with the density $\rho^*$ of Theorem~\ref{thm:equilibrium} for $p=1$ (Remark~\ref{rem:monien}). His
argument assumes existence, checks only stationarity and does not contain an upper bound or a discretisation step.
Theorem~\ref{thm:equilibrium} proves optimality and uniqueness for every $p>0$, and
Sections~\ref{sec:upper}--\ref{sec:lower} supply the passage from the continuous problem to the determinants.

\emph{Averaged arcsine laws.} The measure $\mu^*$ is an average of arcsine laws. A theorem of Kuijlaars and Van
Assche on recurrence coefficients that vary regularly with $n/N$ (\cite{KVA-JAT}, quoted from
\cite[Thm.~1.3]{CCV}) produces measures of exactly this shape as limiting zero distributions. With their
conventions, a family of Jacobi matrices whose diagonal and off-diagonal entries (after rescaling by $N^p$) behave
like $2c\,t^{-p}$ and $c\,t^{-p}$ as $n/N\to t$ has, at $t=1$, the limiting zero distribution $\int_0^1\omega_{[0,4ct^{-p}]}\dd t$. Heuristically, therefore,
$\mu^*$ is the zero distribution that the measures of Theorem~\ref{thm:main} would have if their recurrence
coefficients satisfied $a_n\sim2b_n\sim2c(p)Cn^{-p}$ pointwise; the saturation condition in
Theorem~\ref{thm:equilibrium}(i) is exactly the equation that fixes $c=c(p)$. This heuristic is not used in the
proofs, and we do not prove pointwise asymptotics of the recurrence coefficients: Corollary~\ref{cor:jacobi} is a
statement about double Ces\`aro means only.

\subsection{Method}
By Heine's formula (Lemma~\ref{lem:heine}), $D_n$ is a sum over $n$-subsets $S$ of atoms of
$\prod_{k\in S}w_k$ times a squared Vandermonde product. For $x_k=k^{-p}$ and $S=\{k_1<\dots<k_n\}$, the Vandermonde
product is $n^{-pn(n-1)}\prod_{i\ne j}|g(u_i)-g(u_j)|$ with $g(u)=u^{-p}$ and $u_i=k_i/n$, and the empirical measure
$\frac1n\sum\delta_{u_i}$ has at most one atom in each interval of length $1/n$, a discrete form of the constraint
$\mu\le\Leb$. The upper bound (Section~\ref{sec:upper}) compares point configurations with step densities
(Lemma~\ref{lem:disc}) and removes atoms beyond $Mn$ by a tail estimate; the lower bound
(Section~\ref{sec:lower}) evaluates a single term of the Heine sum at the quantiles of a truncation of $\mu^*_u$. Both
bounds meet at $\sup_\cA\cJ_p$, which Theorem~\ref{thm:equilibrium} identifies. General regularly spaced atoms are
reduced to $x_k=Ck^{-p}$ by a mean-value comparison (Lemma~\ref{lem:spacing}).

The only potential-theoretic input beyond elementary calculus is the positivity of the logarithmic energy of
signed measures of total mass zero (Lemma~\ref{lem:positivity}), which is used only for the optimality and uniqueness
part of Theorem~\ref{thm:equilibrium}(iii). It is a special case of a lemma of Beckermann, Kalyagin, Matos and
Wielonsky~\cite[Lemma~2.1]{BKMW}; we include a short Fourier-analytic proof so that the paper is self-contained
(Remark~\ref{rem:BKMW}).

\subsection{Organisation and notation}
Section~\ref{sec:prelim} collects the facts about Hankel determinants that we use. Section~\ref{sec:equilibrium}
proves Theorem~\ref{thm:equilibrium}. Sections~\ref{sec:upper} and~\ref{sec:lower} prove the upper and lower bounds
for $x_k=k^{-p}$, and Section~\ref{sec:proofs} completes the proofs of Theorem~\ref{thm:main} and
Corollary~\ref{cor:jacobi}. Section~\ref{sec:zeta} treats zeta values, and Section~\ref{sec:numerics} reports
numerical illustrations.

Throughout, $p>0$ is fixed and $g(u):=u^{-p}$ on $(0,\infty)$. Constants denoted $C_M$, $C_U$, \dots\ depend only on
$p$ and on the indicated parameters, and may change from line to line. We use the convention $0\log0=0$.

\section{Hankel determinants of discrete measures}\label{sec:prelim}

In this section $\sigma=\sum_{k\ge1}w_k\delta_{x_k}$ with $w_k\ge0$, $\sum_kw_k<\infty$ and $\sup_k|x_k|<\infty$, so
that all moments are finite. For a finite set $S\subset\N$ write
\[
  w(S):=\prod_{k\in S}w_k,\qquad V_x(S):=\prod_{\substack{k,l\in S\\k<l}}(x_k-x_l)^2 .
\]

\begin{lemma}[Heine's formula]\label{lem:heine}
For every $n\ge1$,
\[
  D_n(\sigma)=\sum_{\substack{S\subset\N\\|S|=n}}w(S)\,V_x(S).
\]
In particular all terms are nonnegative.
\end{lemma}

\begin{proof}
Write $m_{i+j}=\sum_kw_kx_k^ix_k^j$. Expanding the determinant multilinearly in its columns (legitimate because
$\sum_kw_k|x_k|^r\le WX^r$ with $W=\sum_kw_k$ and $X:=\sup_k|x_k|$, so all series converge absolutely),
\[
  D_n=\sum_{k_0,\dots,k_{n-1}\ge1}\Bigl(\prod_{j=0}^{n-1}w_{k_j}x_{k_j}^{\,j}\Bigr)\det\bigl(x_{k_j}^{\,i}\bigr)_{i,j=0}^{n-1}.
\]
The Vandermonde determinant vanishes unless the $k_j$ are distinct. Group the remaining tuples according to the set
$S=\{l_0<\dots<l_{n-1}\}$ they enumerate, $k_j=l_{\pi(j)}$ with $\pi$ a permutation. Then
$\det(x_{k_j}^i)=\operatorname{sgn}(\pi)\det(x_{l_j}^i)$ and
\[
  \sum_\pi\operatorname{sgn}(\pi)\prod_jx_{l_{\pi(j)}}^{\,j}=\det\bigl(x_{l_i}^{\,j}\bigr),
\]
so the contribution of $S$ is $w(S)\det(x_{l_j}^i)^2=w(S)V_x(S)$.
\end{proof}

\begin{lemma}\label{lem:elem}
Let $\sigma,\sigma'$ be measures as above.
\begin{enumerate}
\item (Monotonicity) If $\sigma'=\sum w'_k\delta_{x_k}$ with $w_k\le w'_k$ for all $k$, then
  $D_n(\sigma)\le D_n(\sigma')$.
\item (Translation) $D_n(\sum w_k\delta_{x_k+t})=D_n(\sigma)$ for every $t\in\R$.
\item (Scaling) $D_n(\sum w_k\delta_{\lambda x_k})=\lambda^{n(n-1)}D_n(\sigma)$ for $\lambda>0$, and
  $D_n(\sum\vartheta w_k\delta_{x_k})=\vartheta^nD_n(\sigma)$ for $\vartheta>0$.
\item (Positivity) If $w_k>0$ for infinitely many $k$ with pairwise distinct $x_k$, then $D_n(\sigma)>0$ for all $n$.
\end{enumerate}
\end{lemma}

\begin{proof}
All four statements follow from Lemma~\ref{lem:heine}: its terms are nonnegative and increasing in each $w_k$;
$V_x(S)$ depends only on differences; $V_{\lambda x}(S)=\lambda^{n(n-1)}V_x(S)$ and $(\vartheta w)(S)=\vartheta^nw(S)$; and under the
hypothesis of (iv) some term is positive.
\end{proof}

\begin{lemma}[Regular spacing comparison]\label{lem:spacing}
Let $f,h\in C^1[1,\infty)$ with $f'<0$ and $h'<0$ on $[1,\infty)$, both bounded on $[1,\infty)$, and suppose that
$\rho:=f'/h'$ satisfies $R^{-1}\le\rho\le R$ on $[1,\infty)$ for some $R\ge1$ and $\rho(\xi)\to1$ as $\xi\to\infty$.
Put $\eta(m):=\sup_{\xi\ge m}|\log\rho(\xi)|$. Then for all weights $w_k\ge0$ with $\sum w_k<\infty$ and all $n\ge1$,
\[
  \Bigl|\log D_n\Bigl(\sum_kw_k\delta_{f(k)}\Bigr)-\log D_n\Bigl(\sum_kw_k\delta_{h(k)}\Bigr)\Bigr|
  \le2\sum_{i=1}^{n-1}(n-i)\,\eta(i)=o(n^2),
\]
whenever one (hence both) of the determinants is positive.
\end{lemma}

\begin{proof}
Let $S=\{k_1<\dots<k_n\}$. For $i<j$, Cauchy's mean value theorem gives
$(f(k_i)-f(k_j))/(h(k_i)-h(k_j))=\rho(\xi)$ for some $\xi\in(k_i,k_j)$. Since $\xi>k_i\ge i$, the ratio of the
squared factors lies in $[e^{-2\eta(i)},e^{2\eta(i)}]$. Multiplying over pairs, $V_f(S)/V_h(S)$ lies in
$[e^{-E_n},e^{E_n}]$ with $E_n=2\sum_{i<n}(n-i)\eta(i)$, and the bound follows by comparing the Heine sums term by
term. Finally $\eta$ is bounded by $\log R$ and tends to $0$, so $E_n=o(n^2)$.
\end{proof}

\begin{lemma}[Jacobi coefficients]\label{lem:jacobi}
Suppose $D_n(\sigma)>0$ for all $n$. Let $P_k$ be the monic orthogonal polynomials of $\sigma$ and
$h_k:=\int P_k^2\dd\sigma$. Then $D_n=\prod_{k=0}^{n-1}h_k$. Moreover $xP_k=P_{k+1}+a_{k+1}P_k+b_k^2P_{k-1}$ for
$k\ge1$, with $b_k:=(h_k/h_{k-1})^{1/2}>0$, and
\[
  D_n(\sigma)=m_0^{\,n}\prod_{j=1}^{n-1}b_j^{\,2(n-j)} .
\]
\end{lemma}

The numbers $b_k$ are the off-diagonal entries of the Jacobi matrix of $\sigma$ (the coefficients of the
three-term recurrence for the orthonormal polynomials $P_k/\sqrt{h_k}$).

\begin{proof}
The Hankel matrix $H_n=(m_{i+j})_{i,j<n}$ is the Gram matrix of $1,x,\dots,x^{n-1}$ in $L^2(\sigma)$. Writing
$(P_0,\dots,P_{n-1})^T=L(1,\dots,x^{n-1})^T$ with $L$ unit lower triangular gives
$LH_nL^T=\operatorname{diag}(h_0,\dots,h_{n-1})$, so $D_n=\prod_{k<n}h_k$. The polynomial $xP_k-P_{k+1}$ has degree $\le k$, and it is orthogonal to
$P_l$ for $l\le k-2$; its $P_{k-1}$-coefficient is
\[
  \frac1{h_{k-1}}\int xP_kP_{k-1}\dd\sigma=\frac1{h_{k-1}}\int P_k\,(xP_{k-1})\dd\sigma=\frac{h_k}{h_{k-1}} .
\] Hence $h_k=m_0b_1^2\cdots b_k^2$ and
$D_n=\prod_{k=0}^{n-1}m_0\prod_{j\le k}b_j^2=m_0^n\prod_{j=1}^{n-1}b_j^{2(n-j)}$.
\end{proof}

\begin{lemma}\label{lem:sumlog}
$2\sum_{j=1}^{n-1}(n-j)\log j=n^2\log n-\frac32n^2+O(n\log n)$.
\end{lemma}

\begin{proof}
Let $\phi(x)=(n-x)\log x$ on $[1,n]$. Then $|\phi'(x)|\le\log n+n/x$, so
$\bigl|\sum_{j=1}^{n-1}\phi(j)-\int_1^n\phi\bigr|\le\sum_{j=1}^{n-1}(\log n+n/j)=O(n\log n)$, and
$\int_1^n(n-x)\log x\dd x=\frac{n^2}{2}\log n-\frac34n^2+O(n)$.
\end{proof}

\section{The constrained equilibrium problem}\label{sec:equilibrium}

\subsection{Energies}
For finite positive Borel measures $\mu,\nu$ on $(0,\infty)$ with $\int\log(1+x)\dd\mu<\infty$ and
$\int\log(1+x)\dd\nu<\infty$ put
\[
  \cL(\mu,\nu):=\iint\log|x-y|\dd\mu(x)\dd\nu(y),\qquad\cL(\mu):=\cL(\mu,\mu).
\]
Since $\log|x-y|\le\log(1+x)+\log(1+y)$ for $x,y>0$, the positive part of the integrand is integrable, so
$\cL(\mu,\nu)\in[-\infty,\infty)$ is well defined. Let
\[
  \sigma_p(\dd x):=\tfrac1p\,x^{-1-1/p}\dd x\quad\text{on }(0,\infty),
\]
the image of Lebesgue measure on $(0,\infty)$ under $g(u)=u^{-p}$. For $\tau>0$ let $\cA_\tau$ be the set of positive
Borel measures on $(0,\infty)$ of mass $\tau$ with $\mu\le\Leb$; thus $\cA=\cA_1$. For $\mu\in\cA_\tau$ the image
$\mu^x:=g_*\mu$ satisfies $\mu^x\le\sigma_p$, hence $\int\log(1+x)\dd\mu^x\le\tau\log2+\int_1^\infty\log(1+x)\dd
\sigma_p<\infty$, and we define
\[
  \cJ_p(\mu):=\cL(\mu^x)=\iint\log|g(u)-g(v)|\dd\mu(u)\dd\mu(v)\in[-\infty,\infty),\qquad
  E_p:=\sup_{\mu\in\cA}\cJ_p(\mu).
\]

\begin{lemma}[Scaling]\label{lem:scaling}
For a measure $\mu$ on $(0,\infty)$ and $\tau>0$ let $\mu_\tau(B):=\tau\mu(B/\tau)$. If $\mu$ has mass $m$ and
density $\rho\le1$, then $\mu_\tau$ has mass $\tau m$ and density $\rho(\cdot/\tau)\le1$, and
$\cJ_p(\mu_\tau)=\tau^2(\cJ_p(\mu)-pm^2\log\tau)$. In particular $\mu\mapsto\mu_\tau$ is a bijection
$\cA\to\cA_\tau$ and $\sup_{\cA_\tau}\cJ_p=\tau^2(E_p-p\log\tau)$.
\end{lemma}

\begin{proof}
$\mu_\tau(B)=\tau\int_{B/\tau}\rho(u)\dd u=\int_B\rho(v/\tau)\dd v$. Since $g(\tau u)=\tau^{-p}g(u)$,
\[
  \cJ_p(\mu_\tau)=\tau^2\iint\log|g(\tau u)-g(\tau v)|\dd\mu\dd\mu=\tau^2\bigl(\cJ_p(\mu)-pm^2\log\tau\bigr).
\]
\end{proof}

\subsection{Arcsine potentials}
\begin{lemma}\label{lem:arcsine}
Let $\beta>0$ and $U_\beta(x):=\int\log|x-y|\dd\omega_\beta(y)$. Then $U_\beta(x)=\log(\beta/4)$ for
$x\in[0,\beta]$, and $U_\beta$ is increasing on $[\beta,\infty)$; in particular $U_\beta\ge\log(\beta/4)$ on
$[0,\infty)$. Moreover, for every $x\ge0$,
\begin{equation}\label{eq:absbound}
  \int\bigl|\log|x-y|\bigr|\dd\omega_\beta(y)\le2\log(1+x+\beta)+\bigl|\log(\beta/4)\bigr| .
\end{equation}
\end{lemma}

\begin{proof}
Write $x=\frac\beta2(1-\cos\varphi)$ and $y=\frac\beta2(1-\cos\theta)$ with $\theta,\varphi\in[0,\pi]$. Under
$y=\frac\beta2(1-\cos\theta)$, the measure $\omega_\beta$ is the image of $\dd\theta/\pi$ on $[0,\pi]$, since
$\sqrt{y(\beta-y)}=\frac\beta2\sin\theta$. Now
$\log|x-y|=\log\frac\beta4+\log|2(\cos\theta-\cos\varphi)|$, and
\[
  \log|2(\cos\theta-\cos\varphi)|=\log\Bigl|2\sin\tfrac{\theta+\varphi}{2}\Bigr|
  +\log\Bigl|2\sin\tfrac{\theta-\varphi}{2}\Bigr|
  =-\sum_{k\ge1}\frac2k\cos k\theta\cos k\varphi,
\]
using $\log|2\sin(\alpha/2)|=-\sum_{k\ge1}k^{-1}\cos k\alpha$, the Fourier series of an $L^2(0,2\pi)$ function.
Each of the two series converges in $L^2([0,\pi],\dd\theta)$, hence in $L^1$, so we may integrate termwise; since
$\int_0^\pi\cos k\theta\dd\theta=0$ for $k\ge1$, we get $U_\beta(x)=\log(\beta/4)$ for $x\in[0,\beta]$. For
$x>\beta$, $U_\beta'(x)=\int(x-y)^{-1}\dd\omega_\beta(y)>0$, and $U_\beta$ is continuous at $\beta$ by dominated
convergence (near $y=\beta$ the integrand is dominated by a multiple of $|\log s|\,s^{-1/2}$, $s=\beta-y$). For
\eqref{eq:absbound}, write $|t|=2t^+-t$: since $|x-y|\le x+\beta$ on the support,
$\int|\log|x-y||\dd\omega_\beta=2\int\log^+|x-y|\dd\omega_\beta-U_\beta(x)\le2\log(1+x+\beta)-\log(\beta/4)$.
\end{proof}

\subsection{Positivity of the energy of signed measures}

\begin{lemma}\label{lem:positivity}
Let $\nu$ be a finite signed Borel measure on $\R$ with $\nu(\R)=0$ and
$\iint|\log|x-y||\dd|\nu|(x)\dd|\nu|(y)<\infty$. Then
\[
  -\iint\log|x-y|\dd\nu(x)\dd\nu(y)=\int_0^\infty\frac{|\hat\nu(\xi)|^2}{\xi}\dd\xi\ \ge0,
  \qquad\hat\nu(\xi):=\int e^{i\xi x}\dd\nu(x),
\]
with equality if and only if $\nu=0$.
\end{lemma}

\begin{proof}
For $0<\varepsilon<1<R$ and $a\in\R\setminus\{0\}$ put
$K_{\varepsilon,R}(a):=\int_\varepsilon^R(\cos(a\xi)-e^{-\xi})\,\xi^{-1}\dd\xi$.

\emph{Step 1: a uniform bound.} We use three elementary estimates: (a) for $1\le A<B$,
$|\int_A^B\cos t\,\frac{\dd t}{t}|\le2/A\le2$ (integrate by parts); (b) for $0<A<B\le1$,
$\int_A^B\cos t\,\frac{\dd t}{t}=\log\frac BA-\theta$ with $0\le\theta\le\int_A^B\frac t2\dd t\le\frac{B^2}{4}\le\frac14$,
because $0\le1-\cos t\le t^2/2$;
(c) $\int_\varepsilon^Re^{-\xi}\frac{\dd\xi}{\xi}=\log\frac1\varepsilon+\theta'$ with $|\theta'|\le2$. Substitute
$t=|a|\xi$, so that $\int_\varepsilon^R\cos(a\xi)\frac{\dd\xi}{\xi}=\int_{\varepsilon|a|}^{R|a|}\cos t\frac{\dd t}{t}$.
If $\varepsilon|a|<1<R|a|$, then $K_{\varepsilon,R}(a)=-\log|a|+O(1)$. If $R|a|\le1$, then
$K_{\varepsilon,R}(a)=\log R+O(1)$ and $0<\log R\le\log(1/|a|)$. If $\varepsilon|a|\ge1$, then
$K_{\varepsilon,R}(a)=-\log(1/\varepsilon)+O(1)$ and $0<\log(1/\varepsilon)\le\log|a|$. In all cases
\begin{equation}\label{eq:Kbound}
  |K_{\varepsilon,R}(a)|\le\bigl|\log|a|\bigr|+6\qquad(0<\varepsilon<1<R,\ a\ne0).
\end{equation}

\emph{Step 2: the pointwise limit.} We have
$K_{\varepsilon,R}(a)-K_{\varepsilon,R}(1)=\int_{\varepsilon|a|}^{\varepsilon}\cos t\frac{\dd t}{t}+\int_R^{R|a|}\cos
t\frac{\dd t}{t}$. As $\varepsilon\to0$ the first integral tends to $\log(1/|a|)$, by (b); as $R\to\infty$ the second
tends to $0$, by (a). Also $K_{\varepsilon,R}(1)\to k_0:=\int_0^\infty(\cos\xi-e^{-\xi})\xi^{-1}\dd\xi$, which
converges. Hence $K_{\varepsilon,R}(a)\to k_0-\log|a|$. (In fact $k_0=0$, but we do not need this, because
$\nu(\R)=0$.) All limits as $\varepsilon\to0$, $R\to\infty$ are taken along arbitrary sequences
$\varepsilon_m\downarrow0$, $R_m\uparrow\infty$.

\emph{Step 3.} The integrability hypothesis forces $|\nu|\otimes|\nu|$ to give zero mass to the diagonal $\{x=y\}$.
By \eqref{eq:Kbound} and dominated convergence along $(\varepsilon_m,R_m)$ (with respect to $|\nu|\otimes|\nu|$,
dominating function $|\log|x-y||+6$), and since $\nu(\R)=0$,
\[
  \iint K_{\varepsilon,R}(x-y)\dd\nu\dd\nu\longrightarrow k_0\,\nu(\R)^2-\iint\log|x-y|\dd\nu\dd\nu
  =-\iint\log|x-y|\dd\nu\dd\nu .
\]
On the other hand, the integrand $(\xi,x,y)\mapsto(\cos((x-y)\xi)-e^{-\xi})/\xi$ is bounded on
$[\varepsilon,R]\times\R^2$, so by Fubini, $\iint\cos((x-y)\xi)\dd\nu\dd\nu=|\hat\nu(\xi)|^2$ and $\nu(\R)=0$,
\[
  \iint K_{\varepsilon,R}(x-y)\dd\nu\dd\nu=\int_\varepsilon^R\frac{|\hat\nu(\xi)|^2}{\xi}\dd\xi .
\]
Letting $\varepsilon\to0$, $R\to\infty$ (monotone convergence on the right) gives the identity.

\emph{Equality.} If the integral vanishes, then $\hat\nu=0$ on $(0,\infty)$ by continuity, hence on $\R$, because
$\hat\nu(-\xi)=\overline{\hat\nu(\xi)}$ and $\hat\nu(0)=\nu(\R)=0$. Then $\nu=0$: for the Gaussian
$G_t(x)=(2\pi t)^{-1/2}e^{-x^2/2t}$ we have $G_t(x)=\frac1{2\pi}\int e^{-t\xi^2/2}e^{-i\xi x}\dd\xi$, so by Fubini
$(G_t*\nu)(x)=\frac1{2\pi}\int e^{-t\xi^2/2}e^{-i\xi x}\hat\nu(\xi)\dd\xi=0$ for all $x$ and $t>0$; for bounded
continuous $\psi$, $\int\psi\dd\nu=\lim_{t\to0}\int(G_t*\psi)\dd\nu=\lim_{t\to0}\int\psi\,(G_t*\nu)\dd x=0$ by
dominated convergence, so $\nu=0$.
\end{proof}

\begin{remark}\label{rem:BKMW}
The positivity statement of Lemma~\ref{lem:positivity} (but not the Fourier identity) is a special case of
\cite[Lemma~2.1]{BKMW}. Indeed, under the hypothesis of Lemma~\ref{lem:positivity}, the positive and negative parts
$\nu^\pm$ of $\nu$ have equal mass and finite energy, and they satisfy $\int\log(1+|t|)\dd\nu^\pm(t)<\infty$. For the
last point, fix $y$ outside a $|\nu|$-null set with $\int|\log|x-y||\dd|\nu|(x)<\infty$, and use
$\log(1+|x|)\le\log(1+|y|)+\log 2+\log^+|x-y|$. So $\nu^\pm\in M(\mathbb C)$ in the notation of~\cite{BKMW}. In
\cite{BKMW} the energy is denoted $I(\mu)=\iint\log\frac1{|x-y|}\dd\mu(x)\dd\mu(y)$ (not to be confused with our
constant $I(p)$), $M(\Sigma)$ denotes the Borel measures of finite mass supported
in $\Sigma$ with finite energy and with $\int\log(1+|t|)\dd\mu(t)<\infty$ \cite[(1.12)]{BKMW}, and the statement
reads: ``Lemma~2.1. Let $\mu,\nu\in M(\mathbb C)$ with $\|\mu\|=\|\nu\|$. Then $I(\mu-\nu)\ge0$, and $I(\mu-\nu)=0$ if
and only if $\mu=\nu$.'' The proof given there refers to the literature for both the compactly supported and the
unbounded case. In the application below ($\mu=g_*\mu_0$ with $\mu_0\in\cA$ and $\nu=\mu^*$), both measures are
positive, of mass $1$, of finite energy, and satisfy $\int\log(1+x)<\infty$, so they lie in $M(\mathbb C)$ and
\cite[Lemma~2.1]{BKMW} applies as well. Either reference suffices; we use Lemma~\ref{lem:positivity}.
\end{remark}

\subsection{The explicit maximiser}
Throughout this subsection $c=c(p)$ and $\beta_t:=4ct^{-p}$ for $t\in(0,1]$. Thus $\mu^*$ is the mixture
$\mu^*(B)=\int_0^1\omega_{\beta_t}(B)\dd t$ for Borel sets $B$. Put $U^*(x):=\int\log|x-y|\dd\mu^*(y)$.

\begin{proof}[Proof of Theorem~\ref{thm:equilibrium}(i)]
Fix $x>0$ and let $t_x:=(4c/x)^{1/p}$, so that $x<\beta_t$ if and only if $t<t_x$. Then
\[
  \frac{\dd\mu^*}{\dd x}=\frac1\pi\int_0^{\min(1,t_x)}\frac{\dd t}{\sqrt{x(\beta_t-x)}}
  =\frac{t_x}{\pi x}\int_0^{\min(1,1/t_x)}\frac{\dd s}{\sqrt{s^{-p}-1}},
\]
by the substitution $t=t_xs$, for which $\beta_t=xs^{-p}$. The substitution $v=s^p$ gives, for $0\le V\le1$,
\[
  \int_0^{V^{1/p}}\frac{\dd s}{\sqrt{s^{-p}-1}}=\frac1p\int_0^Vv^{\frac1p-\frac12}(1-v)^{-\frac12}\dd v
  =\frac1pB\bigl(\tfrac1p+\tfrac12,\tfrac12\bigr)\cB_V\bigl(\tfrac1p+\tfrac12,\tfrac12\bigr),
\]
and $B(\frac1p+\frac12,\frac12)=\sqrt\pi\,\Gamma(\frac1p+\frac12)/\Gamma(\frac1p+1)=2I(p)$. With
$V=\min(1,t_x^{-p})=\min(1,x/4c)$ we obtain
\[
  \frac{\dd\mu^*}{\dd x}=\frac{(4c)^{1/p}x^{-1-1/p}}{\pi}\cdot\frac{2I(p)}{p}\cdot\cB_{\min(1,x/4c)}
  \bigl(\tfrac1p+\tfrac12,\tfrac12\bigr)
  =\cB_{\min(1,x/4c)}\bigl(\tfrac1p+\tfrac12,\tfrac12\bigr)\cdot\frac{\dd\sigma_p}{\dd x},
\]
because $(4c)^{1/p}=\pi/(2I(p))$ by the definition of $c=c(p)$. Thus $\mu^*=\sigma_p$ on $[4c,\infty)$ and
$\mu^*<\sigma_p$ (in density) on $(0,4c)$. Pulling back by $x=u^{-p}$, which maps $\sigma_p$ to Lebesgue measure, gives
the density $\rho^*$ of $\mu^*_u$; $\rho^*$ is continuous and positive on $(0,\infty)$, equal to $1$ exactly on
$(0,u_s]$, and $u_s=(4c)^{-1/p}=2I(p)/\pi$. The total mass is $\int_0^1\dd t=1$. Hence $\mu^*_u\in\cA$.
\end{proof}

\begin{proof}[Proof of Theorem~\ref{thm:equilibrium}(ii)]
By \eqref{eq:absbound}, for fixed $x>0$,
$\int_0^1\int|\log|x-y||\dd\omega_{\beta_t}(y)\dd t\le\int_0^1\bigl(2\log(1+x+\beta_t)+|\log(\beta_t/4)|\bigr)\dd t<\infty$,
since $\log\beta_t=\log4c-p\log t$. By Fubini, $U^*(x)=\int_0^1U_{\beta_t}(x)\dd t$. If $x\le4c=\beta_1$, then
$x\le\beta_t$ for all $t\in(0,1]$, and Lemma~\ref{lem:arcsine} gives
$U^*(x)=\int_0^1\log(ct^{-p})\dd t=\log c+p$. For all $x>0$, Lemma~\ref{lem:arcsine} gives
$U_{\beta_t}(x)\ge\log(\beta_t/4)$, hence $U^*(x)\ge\log c+p$.
\end{proof}

\begin{proof}[Proof of Theorem~\ref{thm:equilibrium}(iii)]
\emph{The value.} For $0<s\le t\le1$ we have $\supp\omega_{\beta_t}\subset[0,\beta_t]\subset[0,\beta_s]$, so by
Lemma~\ref{lem:arcsine}
\[
  \iint\log|x-y|\dd\omega_{\beta_s}(x)\dd\omega_{\beta_t}(y)=\int U_{\beta_s}\dd\omega_{\beta_t}=\log\frac{\beta_s}{4}
  =\log c-p\log\min(s,t).
\]
By \eqref{eq:absbound}, $\iint|\log|x-y||\dd\omega_{\beta_s}\dd\omega_{\beta_t}\le2\log(1+\beta_s+\beta_t)+
|\log(\beta_t/4)|$, which is integrable over $(s,t)\in(0,1]^2$. Hence $\cL(\mu^*)$ converges absolutely, and by
Fubini
\begin{align*}
  \cJ_p(\mu^*_u)=\cL(\mu^*)&=\int_0^1\!\!\int_0^1\bigl(\log c-p\log\min(s,t)\bigr)\dd s\dd t\\
  &=\log c-p\int_0^12(1-m)\log m\dd m=\log c+\frac{3p}{2}=\kappa_p .
\end{align*}

\emph{Optimality.} Let $\mu\in\cA$. If $\cJ_p(\mu)=-\infty$ there is nothing to prove, so assume $\cJ_p(\mu)>-\infty$
and let $\mu^x=g_*\mu\le\sigma_p$. We first check that all energies below converge absolutely.
\begin{itemize}
\item $\cL(\mu^x)$: the positive part of the integrand is integrable (Section~3.1), and the negative part is
  integrable because $\cL(\mu^x)>-\infty$.
\item $\cL(\mu^*)$: shown above.
\item $\cL(\mu^x,\mu^*)$: using $|t|=2t^+-t$, $\log^+|x-y|\le\log(1+x)+\log(1+y)$ and part (ii),
  \[
    \int\!\!\int\bigl|\log|x-y|\bigr|\dd\mu^*(y)\dd\mu^x(x)\le\int\Bigl(2\log(1+x)+2\!\int\!\log(1+y)\dd\mu^*(y)
    -(\log c+p)\Bigr)\dd\mu^x(x)<\infty,
  \]
  since $\mu^*\le\sigma_p$ and $\mu^x\le\sigma_p$ have finite logarithmic moments and mass $1$.
\end{itemize}
Let $\nu:=\mu^x-\mu^*$, a finite signed measure on $(0,\infty)$ with $\nu(\R)=0$. Since $|\nu|\le\mu^x+\mu^*$,
the three bullets give $\iint|\log|x-y||\dd|\nu|\dd|\nu|<\infty$. Expanding bilinearly (all terms are finite),
\[
  \cJ_p(\mu)=\cL(\mu^x)=\cL(\mu^*)+2\int U^*\dd\nu+\iint\log|x-y|\dd\nu\dd\nu .
\]
By Lemma~\ref{lem:positivity} the last term is $\le0$, with equality only if $\nu=0$. For the middle term, $\nu$ has
mass $0$, so $\int U^*\dd\nu=\int(U^*-\log c-p)\dd\nu$. By (ii) the integrand vanishes on $(0,4c]$ and is $\ge0$ on
$(4c,\infty)$, where $\nu=\mu^x-\sigma_p\le0$ by (i). Hence $\int U^*\dd\nu\le0$, and
$\cJ_p(\mu)\le\cL(\mu^*)=\kappa_p$, with equality only if $\nu=0$, that is, $\mu=\mu^*_u$.
\end{proof}

\begin{remark}
The only step of the paper that uses Lemma~\ref{lem:positivity} (equivalently \cite[Lemma~2.1]{BKMW}) is the
optimality and uniqueness part of Theorem~\ref{thm:equilibrium}(iii). In the proof of Theorem~\ref{thm:main} it is
needed for the inequality $E_p\le\kappa_p$, which enters the upper bound. The lower bound uses only parts (i) and the
value $\cJ_p(\mu^*_u)=\kappa_p$.
\end{remark}

\begin{remark}\label{rem:monien}
For $p=1$ we have $c=\frac12$ and $\rho^*(u)=\cB_{\min(1,1/2u)}(\frac32,\frac12)$. We checked numerically (mpmath,
30 digits; script \texttt{checks/monien\_check.py}) that this agrees with Monien's density $\rho$ quoted in
Section~\ref{sec:context} at eleven points $x\in[0.1,10^4]$, to within $10^{-30}$.
\end{remark}

\section{Discretisation and the upper bound}\label{sec:upper}

In Sections~\ref{sec:upper} and~\ref{sec:lower} we write, for a finite $S=\{k_1<\dots<k_N\}\subset\N$,
\[
  V(S):=\prod_{i<j}\bigl(k_i^{-p}-k_j^{-p}\bigr)^2,\qquad u_i:=k_i/n .
\]
Since $k^{-p}=n^{-p}g(k/n)$,
\begin{equation}\label{eq:VS}
  \log V(S)=-pN(N-1)\log n+\sum_{i\ne j}\log|g(u_i)-g(u_j)| .
\end{equation}

\begin{lemma}[Discretisation]\label{lem:disc}
Let $M\ge1$. There is $C_M$ such that for all $n\ge2$ and all $S=\{k_1<\dots<k_N\}\subset\N\cap[1,Mn]$ with
$N\le n$,
\[
  \sum_{i\ne j}\log|g(u_i)-g(u_j)|\le n^2\cJ_p(\bar\mu_S)+C_Mn\log n,
\]
where $\bar\mu_S\in\cA_{N/n}$ has density $1$ on $\bigcup_{k\in S}A_k$, $A_k:=[(k-1)/n,k/n]$.
\end{lemma}

\begin{proof}
For cells $A,A'$ write $\langle F\rangle_{A\times A'}:=n^2\iint_{A\times A'}F(u,v)\dd u\dd v$ for the average of
$F(u,v):=\log|g(u)-g(v)|$. Then
$n^2\cJ_p(\bar\mu_S)=\sum_{i\ne j}\langle F\rangle_{A_{k_i}\times A_{k_j}}+\sum_i\langle F\rangle_{A_{k_i}^2}$. Recall
that $|g'(u)|=pu^{-p-1}$ is decreasing.

\emph{Non-adjacent pairs.} Let $k<l$ with $l\ge k+2$. For $u\in A_k$, $v\in A_l$ we have
$g(u)-g(v)\ge g(k/n)-g((l-1)/n)=:D>0$. With $\Delta:=g((l-1)/n)-g(l/n)\le|g'((l-1)/n)|/n$ and
$D\ge\frac{l-1-k}{n}|g'((l-1)/n)|$ (mean value theorem and monotonicity of $|g'|$),
\[
  \log\bigl(g(k/n)-g(l/n)\bigr)=\log(D+\Delta)\le\log D+\frac\Delta D\le\langle F\rangle_{A_k\times A_l}+\frac1{l-k-1}.
\]

\emph{Adjacent pairs.} Let $l=k+1$. The point value is $\log(g(k/n)-g((k+1)/n))\le\log(|g'(k/n)|/n)$. For
$u\in A_k$, $v\in A_{k+1}$, $g(u)-g(v)\ge|g'((k+1)/n)|(v-u)$, so
$\langle F\rangle_{A_k\times A_{k+1}}\ge\log|g'((k+1)/n)|-\log n+\vartheta_0$ with
$\vartheta_0:=\int_0^1\int_1^2\log(y-x)\dd y\dd x=2\log2-\frac32$. Since $|g'(k/n)|/|g'((k+1)/n)|=(1+1/k)^{p+1}\le2^{p+1}$, the
point value exceeds the average by at most $(p+1)\log2-\vartheta_0$.

\emph{Diagonal cells.} For $u,v\in A_k$, $|g(u)-g(v)|\ge|g'(k/n)||u-v|$, so
$\langle F\rangle_{A_k^2}\ge\log|g'(k/n)|-\log n-\frac32$, using $\iint_{[0,1]^2}\log|x-y|=-\frac32$. For $k\le Mn$,
$|g'(k/n)|\ge pM^{-p-1}$.

\emph{Summation.} For fixed $i$, the partners $j>i$ with $k_j\ge k_i+2$ have distinct values $k_j-k_i-1\ge1$, so
$\sum_j1/(k_j-k_i-1)\le1+\log n$; there is at most one adjacent partner. Counting ordered pairs,
\begin{multline*}
  \sum_{i\ne j}\log|g(u_i)-g(u_j)|\le n^2\cJ_p(\bar\mu_S)+2N\bigl(1+\log n+(p+1)\log2+|\vartheta_0|\bigr)\\
  +N\bigl(\log n+\tfrac32+(p+1)\log M+|\log p|\bigr),
\end{multline*}
which is of the stated form since $N\le n$.
\end{proof}

The diagonal estimate is the only place where $S\subset[1,Mn]$ is used: an isolated point $k\gg n$ makes
$\cJ_p(\bar\mu_S)$ very negative while contributing little to the point sum. Atoms beyond $Mn$ are therefore treated
separately in Theorem~\ref{thm:upper}.

\begin{corollary}\label{cor:head}
For $M\ge1$ there is $C'_M$ such that for $n\ge2$ and $S\subset\N\cap[1,Mn]$ with $|S|=N\le n$,
\[
  \log V(S)\le-pN^2\log N+E_pN^2+C'_Mn\log n .
\]
\end{corollary}

\begin{proof}
For $N=0$ both sides are trivial. Otherwise, by \eqref{eq:VS}, Lemma~\ref{lem:disc} and Lemma~\ref{lem:scaling} with
$\tau=N/n$,
\[
  \log V(S)\le-pN(N-1)\log n+N^2\bigl(E_p-p\log(N/n)\bigr)+C_Mn\log n,
\]
and $-pN(N-1)\log n-pN^2\log(N/n)=-pN^2\log N+pN\log n$.
\end{proof}

Note that only the value of the supremum $E_p$ enters here; no information about maximisers is used.

\begin{theorem}[Upper bound]\label{thm:upper}
Let $w_k\ge0$ with $W=\sum_kw_k<\infty$ and $\sigma_0:=\sum_kw_k\delta_{k^{-p}}$. Then
\[
  \limsup_{n\to\infty}\frac1{n^2}\Bigl(\log D_n(\sigma_0)+pn^2\log n\Bigr)\le E_p .
\]
\end{theorem}

\begin{proof}
Fix $M\ge2$. Each $n$-set $S\subset\N$ splits uniquely as $S=H\cup T$ with $H:=S\cap[1,Mn]$ and $T:=S\cap(Mn,\infty)$;
let $m:=|T|$ and $N:=n-m=|H|$. Then $V(S)=V(H)V(T)\prod_{a\in H,b\in T}(a^{-p}-b^{-p})^2$.

\emph{Cross pairs.} For $a\in H$ and $b\in T$, $0<a^{-p}-b^{-p}<a^{-p}$. Since $H$ consists of $N$ distinct positive
integers, $\prod_{a\in H}a\ge N!$, so the cross product is at most $(N!)^{-2pm}$.

\emph{Tail.} All atoms $b^{-p}$ with $b\in T$ lie in the interval $(0,(Mn)^{-p})$; hence
\[
  V(T)\le(Mn)^{-pm(m-1)},\qquad\sum_{|T|=m}w(T)\le W^m .
\]

\emph{Head.} $\sum_{|H|=N}w(H)V(H)\le W^N\sup_HV(H)$, and Corollary~\ref{cor:head} bounds $\sup_HV(H)$.

By Lemma~\ref{lem:heine}, $D_n(\sigma_0)\le\sum_{m=0}^nT_m$, where
\[
  \log T_m\le n\log^+W-2pm\log N!-pN^2\log N+E_pN^2-pm(m-1)\log(Mn)+C'_Mn\log n .
\]
Recall that $E_p=\kappa_p$ is finite (Theorem~\ref{thm:equilibrium}(iii)). Put $t:=m/n$. If $N\ge1$, use $\log N!\ge N\log N-N$, $\log N=\log n+\log(1-t)$ and $\log(Mn)=\log n+\log M$.
The coefficient of $\log n$ is $-p(2mN+N^2+m^2)=-pn^2$, and the remaining terms give
\[
  \log T_m+pn^2\log n-E_pn^2\le n^2\phi_M(t)+O_M(n\log n),
\]
\[
  \phi_M(t):=-p(1-t)(1+t)\log(1-t)+2pt(1-t)-E_p\,t(2-t)-pt^2\log M .
\]
Since $-(1-t)\log(1-t)\le t$ on $[0,1)$, we get $\phi_M(t)\le p(1+t)t+2pt+2|E_p|t-pt^2\log M\le At-pt^2\log M$ with
$A:=4p+2|E_p|$. If $N=0$ ($t=1$), then directly $\log T_n\le n\log^+W-pn(n-1)\log(Mn)+C'_Mn\log n$, and
$-p\log M-E_p\le A-p\log M$, so the same bound holds. The $O_M(n\log n)$ terms are uniform in $m$, and
$\max_{t\in[0,1]}(At-pt^2\log M)\le A^2/(4p\log M)$. Summing the $n+1$ terms,
\[
  \limsup_{n\to\infty}\frac1{n^2}\Bigl(\log D_n(\sigma_0)+pn^2\log n\Bigr)\le E_p+\frac{A^2}{4p\log M}.
\]
Let $M\to\infty$.
\end{proof}

\section{The lower bound}\label{sec:lower}

\begin{lemma}\label{lem:bathtub}
Let $\varphi$ be decreasing on $(0,\infty)$ and integrable near $0$, and let $\lambda\le\Leb$ be a measure on
$(0,\infty)$ of mass $\delta>0$ with $\int|\varphi|\dd\lambda<\infty$. Then $\int\varphi\dd\lambda\le\int_0^\delta
\varphi(u)\dd u$.
\end{lemma}

\begin{proof}
$\int\varphi\dd\lambda-\int_0^\delta\varphi=\int_{(0,\delta]}\varphi\dd(\lambda-\Leb)+\int_{(\delta,\infty)}\varphi\dd
\lambda\le\varphi(\delta)\bigl(\lambda((0,\delta])-\delta\bigr)+\varphi(\delta)\lambda((\delta,\infty))=0$, because
$\lambda-\Leb\le0$ and $\varphi\ge\varphi(\delta)$ on $(0,\delta]$, and $\varphi\le\varphi(\delta)$ on
$(\delta,\infty)$.
\end{proof}

\begin{theorem}[Lower bound]\label{thm:lower}
Let $w_k\ge c_1k^{-s}$ for all $k$, with $c_1>0$, $s\in\R$, and let $\sigma_0=\sum_kw_k\delta_{k^{-p}}$ (with
$\sum w_k<\infty$, so that the moments are finite). Then
\[
  \liminf_{n\to\infty}\frac1{n^2}\Bigl(\log D_n(\sigma_0)+pn^2\log n\Bigr)\ge\kappa_p .
\]
\end{theorem}

\begin{proof}
Fix $\eta>0$.

\emph{Step 1: a compactly supported admissible approximant.} By Theorem~\ref{thm:equilibrium}(i), $\mu^*_u$ has a
continuous positive density $\rho^*\le1$ on $(0,\infty)$. For $U>0$ let $\mu_U:=\mu^*_u|_{(0,U]}$, of mass $m_U$. By
the proof of Theorem~\ref{thm:equilibrium}(iii), $|\log|g(u)-g(v)||$ is $\mu^*_u\otimes\mu^*_u$-integrable, so
by dominated convergence $\cJ_p(\mu_U)\to\cJ_p(\mu^*_u)=\kappa_p$ and $m_U\to1$ as $U\to\infty$. Let
$\theta:=(\mu_U)_{1/m_U}$ (notation of Lemma~\ref{lem:scaling}). Then $\theta$ has mass $1$, density
$\rho^*(um_U)\in[d_U,1]$ on its support $(0,U']$, where $U':=U/m_U$ and $d_U:=\min_{(0,U]}\rho^*\in(0,1]$, and density
$0$ beyond. So $\theta\in\cA$, and by Lemma~\ref{lem:scaling}
\[
  \cJ_p(\theta)=m_U^{-2}\cJ_p(\mu_U)+p\log m_U\ge\kappa_p-\eta
\]
for $U$ large enough, which we fix. Constants $C_U$ below depend on $p$ and $U$ only.

\emph{Step 2: the quantile configuration.} Let $F$ be the distribution function of $\theta$; it is continuous and
strictly increasing on $[0,U']$. For $1\le i\le n$ let $B_i=[l_i,r_i]:=[F^{-1}(\frac{i-1}{n}),F^{-1}(\frac in)]$,
$q_i:=F^{-1}(\frac{i-1/2}{n})$, $k_i:=\lceil nq_i\rceil$ and $u_i:=k_i/n\in[q_i,q_i+\frac1n)$.
\begin{itemize}
\item Since $\theta\le\Leb$ and $\theta(B_i)=\frac1n$: $|B_i|\ge\frac1n$, so $l_i\ge\frac{i-1}{n}$ and
  $r_i\ge\frac in$; also $r_i-q_i\ge\frac1{2n}$ and $q_i-l_i\ge\frac1{2n}$. Hence $q_{i+1}-q_i\ge\frac1n$, so the $k_i$
  are distinct, $k_1\ge1$, and $u_j-u_i\ge\frac1n$ for $i<j$. Moreover, for $i\le n-1$, $u_i<q_i+\frac1n\le r_i+\frac1{2n}\le r_{i+1}$.
\item Since the density is $\ge d_U$ on $(0,U']$: $|B_i|\le\frac1{d_Un}$, so $r_i\le\frac{i}{d_Un}$, and
  $u_j\le r_j+\frac1{2n}\le\frac{j+1}{d_Un}$.
\item $S^*:=\{k_1<\dots<k_n\}\subset[1,U'n+1]$.
\end{itemize}

\emph{Step 3: energy comparison.} Write $F(u,v)=\log|g(u)-g(v)|$ and, for blocks $B,B'$, put
$\langle F\rangle_{B\times B'}:=n^2\iint_{B\times B'}F\dd\theta\dd\theta$; this is an average, since
$\theta(B_i)=\frac1n$.
Then
\[
  n^2\cJ_p(\theta)=\sum_{i\ne j}\langle F\rangle_{B_i\times B_j}+\sum_i\langle F\rangle_{B_i^2}.
\]
For $i<j$,
$u\in B_i$ and $v\in B_j$ we have $u\le v$ and
\[
  F(u,v)=-p\log u+h(u/v),\qquad h(r):=\log(1-r^p),
\]
with $h$ decreasing on $(0,1)$. We show that $F(u_i,u_j)\ge\langle F\rangle_{B_i\times B_j}-\mathrm{err}_{ij}$
with $\sum_{i<j}\mathrm{err}_{ij}=O_U(n\log n)$. We use repeatedly that for $0<r<1$,
\begin{equation}\label{eq:rp}
  \min(1,p)(1-r)\le1-r^p\le\max(1,p)(1-r).
\end{equation}

(a) \emph{The term $-p\log u$.} The point value minus the average is $-p\langle\log(u_i/u)\rangle_{B_i}$, and
$u_i\le r_{i+1}$, $u\ge l_i$ on $B_i$. For $i\ge2$, $\log(r_{i+1}/l_i)\le(r_{i+1}-l_i)/l_i\le\frac{2/(d_Un)}{(i-1)/n}$,
so the loss is at most $2p/(d_U(i-1))$. For $i=1$, Lemma~\ref{lem:bathtub} (with $\varphi=\log(1/u)$ and
$\lambda=\theta|_{B_1}$, $\delta=1/n$) gives $\langle\log(1/u)\rangle_{B_1}\le n\int_0^{1/n}\log\frac1u\dd u=\log n+1$,
and $\log u_1\le\log r_2\le\log\frac2{d_Un}$; so the loss is at most $p(\log\frac2{d_U}+1)$. Summing over $j>i$ and
then over $i$ gives $O_U(n\log n)$.

(b) \emph{The term $h$, for $j\ge i+3$.} On $B_i\times B_j$, $u/v\ge a:=l_i/r_j$, so the average of $h(u/v)$ is at most
$h(a)$. At the point, $u_i/u_j\le b:=r_{i+1}/l_j$, and $b<1$ because $l_j=r_{j-1}\ge r_{i+2}>r_{i+1}$; so
$h(u_i/u_j)\ge h(b)$. The loss is at most $h(a)-h(b)=\log\frac{1-a^p}{1-b^p}\le\frac{b^p-a^p}{1-b^p}$. By
\eqref{eq:rp}, $b^p-a^p=b^p(1-(a/b)^p)\le\max(1,p)(1-a/b)$.
\begin{itemize}
\item If $b\le\frac12$: $1-b^p\ge1-2^{-p}$, and
  $1-\frac ab\le\bigl(1-\frac{l_i}{r_{i+1}}\bigr)+\bigl(1-\frac{l_j}{r_j}\bigr)\le\frac{2/(d_Un)}{(i+1)/n}+
  \frac{1/(d_Un)}{j/n}\le\frac{3}{d_U(i+1)}$. The loss is at most $C_U/(i+1)$.
\item If $b>\frac12$: by \eqref{eq:rp}, $1-b^p\ge\min(1,p)(1-b)=\min(1,p)\frac{l_j-r_{i+1}}{l_j}\ge
  \min(1,p)\frac{j-i-2}{nl_j}$, since the $j-i-2$ blocks $B_{i+2},\dots,B_{j-1}$ each have length $\ge\frac1n$. Also
  $1-\frac ab=\frac{b-a}{b}\le2(b-a)$ and $b-a\le\frac{r_{i+1}-l_i}{l_j}+\frac{l_i}{r_j}\cdot\frac{r_j-l_j}{l_j}\le
  \frac{3}{d_Unl_j}$. The loss is at most $C_U/(j-i-2)$.
\end{itemize}
Summing, $\sum_{j\ge i+3}\bigl(C_U/(i+1)+C_U/(j-i-2)\bigr)=O_U(n\log n)$.

(c) \emph{The term $h$, for $j\in\{i+1,i+2\}$.} The average of $h$ is $\le0$ since $h\le0$. At the point, by
\eqref{eq:rp}, $1-(u_i/u_j)^p\ge\min(1,p)\frac{u_j-u_i}{u_j}\ge\min(1,p)\frac{1/n}{(j+1)/(d_Un)}=
\frac{\min(1,p)d_U}{j+1}$. So the loss is at most $\log(j+1)+C_U$, and there are at most $2n$ such pairs: $O_U(n\log n)$.

(d) \emph{Diagonal blocks.} $|g(u)-g(v)|\le g(\min(u,v))$, so $\langle F\rangle_{B_i^2}\le
p\langle\log(1/\min(u,v))\rangle_{B_i^2}$. For $i\ge2$, $\min(u,v)\ge l_i\ge\frac{i-1}{n}$, giving
$\le p\log\frac{n}{i-1}$. For $i=1$ and $n\ge1/d_U$ we have $B_1\subset(0,1]$, so
$\log\frac1{\min(u,v)}\le\log\frac1u+\log\frac1v$, and Lemma~\ref{lem:bathtub} gives $\le2p(\log n+1)$. In total
$\sum_i\langle F\rangle_{B_i^2}\le p\bigl((n-1)\log n-\log(n-1)!\bigr)+2p(\log n+1)=O(n)$.

Combining (a)--(d),
\begin{align*}
  \sum_{i\ne j}\log|g(u_i)-g(u_j)|=2\sum_{i<j}F(u_i,u_j)
  &\ge n^2\cJ_p(\theta)-\sum_i\langle F\rangle_{B_i^2}-O_U(n\log n)\\
  &\ge n^2\cJ_p(\theta)-O_U(n\log n).
\end{align*}

\emph{Step 4: conclusion.} By Lemma~\ref{lem:heine}, $D_n(\sigma_0)\ge w(S^*)V(S^*)$. Here
$\log w(S^*)\ge n\log c_1-|s|\,n\log(U'n+1)$, and by \eqref{eq:VS} and Step~3,
\[
  \log D_n(\sigma_0)\ge-pn(n-1)\log n+n^2(\kappa_p-\eta)-O_U(n\log n).
\]
Let $n\to\infty$ and then $\eta\to0$.
\end{proof}

The lower bound uses Theorem~\ref{thm:equilibrium}(i) and the value $\cJ_p(\mu^*_u)=\kappa_p$, but neither
Lemma~\ref{lem:positivity} nor any optimality statement.

\section{Proof of Theorem~\ref{thm:main} and Corollary~\ref{cor:jacobi}}\label{sec:proofs}

\begin{proof}[Proof of Theorem~\ref{thm:main}]
\emph{Pure powers.} Let $\sigma_0=\sum_kw_k\delta_{k^{-p}}$ with weights satisfying \eqref{eq:Hw}. By
Theorems~\ref{thm:upper}, \ref{thm:lower} and~\ref{thm:equilibrium}(iii) ($E_p=\kappa_p$),
$\log D_n(\sigma_0)=-pn^2\log n+\kappa_pn^2+o(n^2)$. By Lemma~\ref{lem:elem}(iii), the measure
$\sum_kw_k\delta_{Ck^{-p}}$ has $\log D_n$ larger by $n(n-1)\log C=n^2\log C+O(n)$.

\emph{General atoms.} Let $f$ satisfy \eqref{eq:Ha} and put $h(\xi):=C\xi^{-p}$. Since $f'$ is continuous,
$\xi^{p+1}|f'(\xi)|$ is bounded below by a positive constant, and $\xi^{p+1}f'(\xi)\to-pC<0$, we have $f'<0$ on
$[1,\infty)$; $f$ is bounded because $|f'(\xi)|\le C'\xi^{-p-1}$ is integrable. The ratio
$\rho=f'/h'=\xi^{p+1}f'(\xi)/(-pC)$ is bounded above and below by positive constants and tends to $1$. All
determinants involved are positive by Lemma~\ref{lem:elem}(iv), since $w_k>0$ for all $k$ and $f$, $h$ are injective.
Lemma~\ref{lem:spacing} gives $\log D_n(\sigma)=\log D_n(\sum_kw_k\delta_{Ck^{-p}})+o(n^2)$.
\end{proof}

\begin{proof}[Proof of Corollary~\ref{cor:jacobi}]
By Lemma~\ref{lem:jacobi}, $\log D_n=n\log m_0+2\sum_{j<n}(n-j)\log b_j$, so
\[
  \frac{2}{n^2}\sum_{j=1}^{n-1}(n-j)\log(j^pb_j)=\frac1{n^2}\Bigl(\log D_n-n\log m_0+2p\sum_{j=1}^{n-1}(n-j)\log j\Bigr).
\]
By Theorem~\ref{thm:main} and Lemma~\ref{lem:sumlog}, the bracket equals
$-pn^2\log n+(\kappa_p+\log C)n^2+pn^2\log n-\frac{3p}{2}n^2+o(n^2)=\log\bigl(c(p)C\bigr)n^2+o(n^2)$.
\end{proof}

\begin{remark}
Theorem~\ref{thm:main} gives the error $o(n^2)$ only. For pure powers the proof gives no rate, because the rate of
$\cJ_p(\theta)\to\kappa_p$ in Step~1 of Theorem~\ref{thm:lower} and the factor $A^2/(4p\log M)$ in
Theorem~\ref{thm:upper} are not quantified. The numerics of Section~\ref{sec:numerics} are consistent with an error
$O(n\log n)$ for pure powers.
\end{remark}

\section{Hankel determinants of zeta values}\label{sec:zeta}

\begin{proof}[Proof of Corollary~\ref{cor:zeta}]
Since $2a+b>1$, every entry $\zeta(a(i+j)+b)$ with $i,j\ge1$ is a convergent Dirichlet series, and
\[
  \zeta\bigl(a(i+j)+b\bigr)=\sum_{k\ge1}k^{-(2a+b)}\bigl(k^{-a}\bigr)^{(i-1)+(j-1)} .
\]
Re-indexing $i,j\mapsto i-1,j-1$, the determinant is $D_n(\sigma)$ for $\sigma=\sum_kk^{-(2a+b)}\delta_{k^{-a}}$.
Apply Theorem~\ref{thm:main} with $p=a$, $f(\xi)=\xi^{-a}$ (so $C=1$ and $\xi^{a+1}f'(\xi)\equiv-a$), and
$w_k=k^{-(2a+b)}$, for which $\sum_kw_k<\infty$ and $w_k\ge k^{-s}$ with $s=2a+b$. The special cases follow from
$c(1)=\frac12$ and $c(2)=\frac{\pi^2}{16}$.
\end{proof}

For example, $\frac{3a}{2}+\log c(a)$ equals $0.806853$, $2.516871$, $4.127795$, $5.696954$ (rounded to six decimals)
for $a=1,2,3,4$ (Table~\ref{tab:constants}).

\begin{remark}[$a=1$]
Monien's $H_n^{(r)}[\zeta]=\det_{1\le i,j\le n}\zeta(i+j+r)$ is the case $a=1$, $b=r$ of Corollary~\ref{cor:zeta}, so
$\log H_n^{(r)}[\zeta]=-n^2(\log2n-\frac32)+o(n^2)$ for every real $r>-1$. This is Monien's weak asymptotic form for $r=0,1$
\cite{Monien}, read with an $o(n^2)$ error term, which is how the expansion from Zagier's numerics quoted there
reads it. Monien's ``$\sim$'' by itself asserts only the leading term.
\end{remark}

\begin{remark}[$a=2$: comparison with closed forms]\label{rem:bs}
Baricz and \v{S}tampach~\cite{BariczStampach} evaluate Hankel determinants of the Rayleigh functions
$\sigma_{2m}(\nu)$. (The statements and their numbering below are as recorded in our literature notes on~\cite{BariczStampach} and remain to be re-checked against the published version; the closed forms themselves are confirmed independently below.)
Their Corollary~3 reads ``For $n\in\N$, $\nu\notin-\N$, and $\ell\in\{0,1\}$, one has $\det H_n^{(\ell)}(\nu)=
2^{-2n(n+\ell)}\prod_{k=1}^{2n+\ell-1}(\nu+k)^{k-2n-\ell}$'', where $H_n^{(\ell)}(\nu)$ has $(i,j)$ entry
$\sigma_{2(i+j)+2\ell-2}(\nu)$; their Remark~5 gives $\det H_n^{(2)}(\nu)=2^{-2n(n+2)}(n+1)(n+\nu+1)\prod_{k=1}^{2n+1}
(\nu+k)^{k-2n-2}$; and their (22) reads $\sigma_{2n}(1/2)=\zeta(2n)/\pi^{2n}$. Taking $\nu=\frac12$ and rescaling rows
and columns by powers of $\pi$ yields closed forms for $\det_{1\le i,j\le n}\zeta(2(i+j)+b)$ with $b=2\ell-2\in
\{-2,0,2\}$. For example, for $b=0$,
\[
  \log\det_{1\le i,j\le n}\zeta\bigl(2(i+j)\bigr)=2n(n+1)\log\frac\pi2+\sum_{k=1}^{2n}(k-2n-1)\log\bigl(k+\tfrac12\bigr).
\]
We verified these three closed forms against direct 60-digit determinants for $n\le8$ (agreement to at least
$10^{-38}$; script \texttt{checks/bs\_check.py}). The asymptotics of the closed form are elementary: with
$m=2n+1$, $\sum_{k<m}(m-k)\log(k+\frac12)=\frac{m^2}{2}\log m-\frac34m^2+O(m\log m)$ (as in Lemma~\ref{lem:sumlog}),
which gives
\[
  \log\det_{1\le i,j\le n}\zeta\bigl(2(i+j)\bigr)=-2n^2\log n+\bigl(3+\log\tfrac{\pi^2}{16}\bigr)n^2+O(n\log n),
\]
in agreement with Corollary~\ref{cor:zeta} for $a=2$, $b=0$, and with a better error term. Numerically, for
$\ell=0,1,2$ the quantity $(\log\det+2n^2\log n)/n^2-\kappa_2$ equals $-1.33\cdot10^{-4}$, $-3.48\cdot10^{-4}$ and
$-5.63\cdot10^{-4}$ at $n=10^5$, and $(n/\log n)$ times it lies between $-1.15$ and $-4.89$ for
$n\in\{10^2,10^3,10^4,10^5\}$. The closed forms for $\ell=0$ and $\ell=2$ differ from the case $\ell=1$ only by
the power of $2$ ($2^{-2n(n+\ell)}$), the power of $\pi$ (exponent $2\sum_{i\le n}(2i+\ell-1)$), the range $2n+\ell-1$
of the product (with $m=2n+\ell$ in place of $2n+1$), and, for $\ell=2$, the factor $(n+1)(n+\frac32)$. These changes
alter the right-hand side only by $O(n\log n)$. Hence the same computation gives
$-2n^2\log n+(3+\log\frac{\pi^2}{16})n^2+O(n\log n)$ for $b=-2$ and $b=2$ as well. Thus for $a=2$,
$b\in\{-2,0,2\}$, Corollary~\ref{cor:zeta} (with error $O(n\log n)$) also follows from \cite{BariczStampach}; the
asymptotic statement is not made there.
\end{remark}

\section{Numerical illustrations}\label{sec:numerics}

The computations below are floating-point (not certified) and serve only as illustrations. They were produced with
scripts that will be made available as supplementary material (repository URL and commit hash: to be supplied
before submission). They were re-run for this draft, and the outputs reproduce the stored ones exactly.

\emph{Constants and the maximiser.} Table~\ref{tab:constants} lists $c(p)$, $u_s$ and $\kappa_p$, together with a
Monte Carlo estimate of $\cL(\mu^*)$ obtained by sampling $\mu^*$ directly (independently of the closed-form value).
The density ratio $\dd\mu^*/\dd\sigma_p$ evaluated by quadrature equals $1$ for $x\ge4c$ to four digits (one
exception at $p=0.5$, $x=4000c$, where the quadrature returns $0.9862$; the closed form of
Theorem~\ref{thm:equilibrium}(i) gives exactly $1$ there) and is strictly below $1$ for $x<4c$.

\begin{table}[ht]
\centering
\begin{tabular}{rcccc}
\toprule
$p$ & $c(p)$ & $u_s=2I(p)/\pi$ & $\kappa_p$ & Monte Carlo $\cL(\mu^*)$\\
\midrule
0.5 & 0.408248 & 0.3750 & $-0.145880$ & $-0.1442\pm0.0021$\\
1   & 0.500000 & 0.5000 & $0.806853$ & $0.8054\pm0.0027$\\
1.5 & 0.566272 & 0.5798 & $1.681320$ & $1.6818\pm0.0034$\\
2   & 0.616850 & 0.6366 & $2.516871$ & $2.5208\pm0.0042$\\
3   & 0.689213 & 0.7132 & $4.127795$ & $4.1281\pm0.0058$\\
4   & 0.738565 & 0.7628 & $5.696954$ & $5.6884\pm0.0075$\\
6   & 0.801586 & 0.8235 & $8.778837$ & $8.7781\pm0.0109$\\
\bottomrule
\end{tabular}
\caption{Constants of Theorem~\ref{thm:equilibrium}. Source: supplementary script \texttt{check\_energy}.}
\label{tab:constants}
\end{table}

\emph{The averaged law.} For $x_k=k^{-p}$ and $w_k=k^{-s}$ truncated to $k\le K$, the recurrence coefficients $b_j$
were computed in extended precision, and
$S_n:=\frac{2}{n^2}\sum_{j<n}(n-j)\log(j^pb_j)-\log c(p)$ is shown in Table~\ref{tab:averaged}; by
Corollary~\ref{cor:jacobi}, $S_n\to0$. Doubling $K$ changes the tabulated values by at most $5.4\cdot10^{-4}$
($p=1.5$; this is about $20\%$ of the entry at $n=1600$, which is therefore only indicative),
$2.8\cdot10^{-5}$ ($p=3$), $1.4\cdot10^{-6}$ ($p=4$) and $7\cdot10^{-8}$ ($p=6$). For $p=0.5$ the atoms
decay too slowly for the truncation used ($K$-doubling changes $S_n$ by $10^{-2}$), so that case is not illustrated.
For the zeta weights of Corollary~\ref{cor:zeta} with $(a,b)=(3,1)$ and $(4,1)$ the same statistic is
$-0.5544,-0.2181,-0.0750$ and $-0.7452,-0.2922,-0.1004$ at $n=25,100,400$. The identification of
$\det\zeta(a(i+j)+b)$ with $D_n$ of $\sum_kk^{-(2a+b)}\delta_{k^{-a}}$ was checked to 12 digits for
$(a,b)\in\{(3,1),(4,1),(3,2)\}$ and $n\in\{4,8,12\}$.

\begin{table}[ht]
\centering
\begin{tabular}{rrrrrrrrr}
\toprule
$p$ & $s$ & $K$ & $n=25$ & $50$ & $100$ & $200$ & $400$ & $1600$\\
\midrule
1.5 & 1.5 & 50000 & $-0.03377$ & $-0.02370$ & $-0.01531$ & $-0.00942$ & $-0.00567$ & $-0.00267$\\
3   & 3   & 20000 & $-0.17266$ & $-0.11244$ & $-0.06962$ & $-0.04161$ & $-0.02423$ & $-0.00782$\\
4   & 4   & 6000  & $-0.27151$ & $-0.17494$ & $-0.10757$ & $-0.06399$ & $-0.03714$ & --\\
6   & 6   & 2000  & $-0.46811$ & $-0.29944$ & $-0.18324$ & $-0.10863$ & -- & --\\
\bottomrule
\end{tabular}
\caption{The averaged statistic $S_n$ (supplementary scripts \texttt{check\_averaged}, \texttt{check\_hz}).}
\label{tab:averaged}
\end{table}

\bibliographystyle{amsplain}
\bibliography{refs}

\end{document}
